\subsection{Main comparison}
Table~\ref{tab:main} gives the primary result: the fold-averaged AUROC, AUPRC and Brier score of all five systems, mean $\pm$ standard deviation over \rhval{const/n_seeds} seeds, each seed being a different \rhval{const/n_folds}-fold patient-level split. Gradient boosting on summary features is best on AUROC ( \rhval{main/gbdt-summary/physionet2012_mortality/auroc/mean:4}). It is also best on AUPRC and tied with logistic regression on Brier score (paired test below). Logistic regression is second on AUROC (\rhval{main/lr-summary/physionet2012_mortality/auroc/mean:4}), and the three recurrent models are lower and close to each other: forward-fill GRU \rhval{main/gru-forward-fill/physionet2012_mortality/auroc/mean:4}, mask+delta GRU \rhval{main/gru-mask+delta/physionet2012_mortality/auroc/mean:4}, GRU-D \rhval{main/gru-d/physionet2012_mortality/auroc/mean:4}. On AUPRC the picture is the same for GBDT (\rhval{main/gbdt-summary/physionet2012_mortality/auprc/mean:4}) but the gap between LR (\rhval{main/lr-summary/physionet2012_mortality/auprc/mean:4}) and the GRUs shrinks to about the seed noise, and Brier scores of all systems lie within a narrow band (from \rhval{main/gbdt-summary/physionet2012_mortality/brier/mean:4} for GBDT to \rhval{main/gru-forward-fill/physionet2012_mortality/brier/mean:4} for the forward-fill GRU). Figure~\ref{fig:seed} shows the per-seed values: GBDT is above every recurrent model in every seed.

\begin{table*}[t]\centering\caption{Main results on PhysioNet 2012 set-a, mean $\pm$ sd over seeds of the fold-averaged metric (5-fold patient-level CV per seed). Higher is better for AUROC and AUPRC, lower for Brier. Bold: best; underline: second best (on Brier score GBDT and logistic regression are tied, paired $p=$\rhval{analysis/paired-analysis-v2/physionet2012_mortality/p_gbdt_lr_brier/mean:2}).}\label{tab:main}
\begin{tabular}{lcccc}
\toprule
Method & auroc $\uparrow$ & auprc $\uparrow$ & brier $\downarrow$ & $n$ \\
\midrule
LR (summary) & \underline{0.8429 $\pm$ 0.0034} & \underline{0.4936 $\pm$ 0.0054} & \underline{0.0926 $\pm$ 0.0006} & 5 \\
GBDT (summary) & \textbf{0.8576 $\pm$ 0.0048} & \textbf{0.5072 $\pm$ 0.0112} & \textbf{0.0925 $\pm$ 0.0013} & 5 \\
GRU (forward-fill) & 0.8323 $\pm$ 0.0037 & 0.4837 $\pm$ 0.0096 & 0.0944 $\pm$ 0.0010 & 5 \\
GRU (mask+delta) & 0.8329 $\pm$ 0.0023 & 0.4884 $\pm$ 0.0064 & 0.0941 $\pm$ 0.0004 & 5 \\
\textbf{GRU-D} & 0.8301 $\pm$ 0.0021 & 0.4896 $\pm$ 0.0075 & 0.0942 $\pm$ 0.0009 & 5 \\
\bottomrule
\end{tabular}
\end{table*}

\begin{figure}[t]\centering\includegraphics[width=0.92\columnwidth]{seed_auroc.pdf}\caption{AUROC per seed (dots; grey lines connect the same seed/fold split across systems; black bar = mean) at the full 48 h horizon.}\label{fig:seed}\end{figure}

\subsection{Hypothesis by hypothesis}
All tests are paired over seeds (the same fold split per seed for every system), two-sided paired $t$-tests on the fold-averaged metric with $n=$\,\rhval{const/n_seeds}; they are unadjusted for multiplicity, and with five pairs per test the smallest attainable two-sided $p$ is limited, so we also report the number of seeds in which the first system wins (Table~\ref{tab:paired}).

\begin{table*}[t]\centering\caption{Paired differences at 48 h (first system minus second; wins = seeds with the first system better, of \rhval{const/n_seeds}). $p$: paired $t$-test on AUROC; $p_{\mathrm{Br}}$: on Brier score (negative $\Delta$Brier favours the first system).}\label{tab:paired}
\begin{tabular}{llcccccc}\toprule
H & Pair (48 h) & $\Delta$AUROC & $p$ & wins & $\Delta$AUPRC & $\Delta$Brier & $p_{\mathrm{Br}}$\\\midrule
H1 & GRU-D $-$ GBDT & \rhval{analysis/paired-analysis-v2/physionet2012_mortality/d_grud_gbdt_auroc/mean:4} & \rhval{analysis/paired-analysis-v2/physionet2012_mortality/p_grud_gbdt_auroc/mean:4} & \rhval{analysis/paired-analysis-v2/physionet2012_mortality/w_grud_gbdt_auroc/mean:0} & \rhval{analysis/paired-analysis-v2/physionet2012_mortality/d_grud_gbdt_auprc/mean:4} & \rhval{analysis/paired-analysis-v2/physionet2012_mortality/d_grud_gbdt_brier/mean:4} & \rhval{analysis/paired-analysis-v2/physionet2012_mortality/p_grud_gbdt_brier/mean:4} \\
 & GRU-D $-$ LR & \rhval{analysis/paired-analysis-v2/physionet2012_mortality/d_grud_lr_auroc/mean:4} & \rhval{analysis/paired-analysis-v2/physionet2012_mortality/p_grud_lr_auroc/mean:4} & \rhval{analysis/paired-analysis-v2/physionet2012_mortality/w_grud_lr_auroc/mean:0} & \rhval{analysis/paired-analysis-v2/physionet2012_mortality/d_grud_lr_auprc/mean:4} & \rhval{analysis/paired-analysis-v2/physionet2012_mortality/d_grud_lr_brier/mean:4} & \rhval{analysis/paired-analysis-v2/physionet2012_mortality/p_grud_lr_brier/mean:4} \\
 & GBDT $-$ LR & \rhval{analysis/paired-analysis-v2/physionet2012_mortality/d_gbdt_lr_auroc/mean:4} & \rhval{analysis/paired-analysis-v2/physionet2012_mortality/p_gbdt_lr_auroc/mean:4} & \rhval{analysis/paired-analysis-v2/physionet2012_mortality/w_gbdt_lr_auroc/mean:0} & \rhval{analysis/paired-analysis-v2/physionet2012_mortality/d_gbdt_lr_auprc/mean:4} & \rhval{analysis/paired-analysis-v2/physionet2012_mortality/d_gbdt_lr_brier/mean:4} & \rhval{analysis/paired-analysis-v2/physionet2012_mortality/p_gbdt_lr_brier/mean:4} \\
H2 & GRU(m+d) $-$ GRU(ffill) & \rhval{analysis/paired-analysis-v2/physionet2012_mortality/d_md_ffill_auroc/mean:4} & \rhval{analysis/paired-analysis-v2/physionet2012_mortality/p_md_ffill_auroc/mean:4} & \rhval{analysis/paired-analysis-v2/physionet2012_mortality/w_md_ffill_auroc/mean:0} & \rhval{analysis/paired-analysis-v2/physionet2012_mortality/d_md_ffill_auprc/mean:4} & \rhval{analysis/paired-analysis-v2/physionet2012_mortality/d_md_ffill_brier/mean:4} & \rhval{analysis/paired-analysis-v2/physionet2012_mortality/p_md_ffill_brier/mean:4} \\
 & GRU-D $-$ GRU(ffill) & \rhval{analysis/paired-analysis-v2/physionet2012_mortality/d_grud_ffill_auroc/mean:4} & \rhval{analysis/paired-analysis-v2/physionet2012_mortality/p_grud_ffill_auroc/mean:4} & \rhval{analysis/paired-analysis-v2/physionet2012_mortality/w_grud_ffill_auroc/mean:0} & \rhval{analysis/paired-analysis-v2/physionet2012_mortality/d_grud_ffill_auprc/mean:4} & \rhval{analysis/paired-analysis-v2/physionet2012_mortality/d_grud_ffill_brier/mean:4} & \rhval{analysis/paired-analysis-v2/physionet2012_mortality/p_grud_ffill_brier/mean:4} \\
H3 & GRU-D $-$ GRU(m+d) & \rhval{analysis/paired-analysis-v2/physionet2012_mortality/d_grud_md_auroc/mean:4} & \rhval{analysis/paired-analysis-v2/physionet2012_mortality/p_grud_md_auroc/mean:4} & \rhval{analysis/paired-analysis-v2/physionet2012_mortality/w_grud_md_auroc/mean:0} & \rhval{analysis/paired-analysis-v2/physionet2012_mortality/d_grud_md_auprc/mean:4} & \rhval{analysis/paired-analysis-v2/physionet2012_mortality/d_grud_md_brier/mean:4} & \rhval{analysis/paired-analysis-v2/physionet2012_mortality/p_grud_md_brier/mean:4} \\
 & GRU(ffill) $-$ LR & \rhval{analysis/paired-analysis-v2/physionet2012_mortality/d_ffill_lr_auroc/mean:4} & \rhval{analysis/paired-analysis-v2/physionet2012_mortality/p_ffill_lr_auroc/mean:4} & \rhval{analysis/paired-analysis-v2/physionet2012_mortality/w_ffill_lr_auroc/mean:0} & \rhval{analysis/paired-analysis-v2/physionet2012_mortality/d_ffill_lr_auprc/mean:4} & \rhval{analysis/paired-analysis-v2/physionet2012_mortality/d_ffill_lr_brier/mean:4} & \rhval{analysis/paired-analysis-v2/physionet2012_mortality/p_ffill_lr_brier/mean:4} \\
\bottomrule\end{tabular}
\end{table*}

\textbf{H1: GRU-D does not beat gradient boosting on summaries (supported).} GRU-D is lower than GBDT by \rhval{analysis/paired-analysis-v2/physionet2012_mortality/d_grud_gbdt_auroc/mean:4} AUROC (paired $p=$\rhval{analysis/paired-analysis-v2/physionet2012_mortality/p_grud_gbdt_auroc/mean:4}, GRU-D is better in \rhval{analysis/paired-analysis-v2/physionet2012_mortality/w_grud_gbdt_auroc/mean:0} of the five seeds), and by \rhval{analysis/paired-analysis-v2/physionet2012_mortality/d_grud_gbdt_auprc/mean:4} AUPRC. The registered refutation criterion (GRU-D better by more than \rhval{const/margin_auroc:2} with $p<\rhval{const/alpha:2}$) is not met; the observed difference is in the opposite direction. GRU-D is also below logistic regression on AUROC (\rhval{analysis/paired-analysis-v2/physionet2012_mortality/d_grud_lr_auroc/mean:4}, $p=$\rhval{analysis/paired-analysis-v2/physionet2012_mortality/p_grud_lr_auroc/mean:4}), although the AUPRC difference to LR (\rhval{analysis/paired-analysis-v2/physionet2012_mortality/d_grud_lr_auprc/mean:4}, $p=$\rhval{analysis/paired-analysis-v2/physionet2012_mortality/p_grud_lr_auprc/mean:4}) is not distinguishable from zero with five seeds. Gradient boosting beats LR (\rhval{analysis/paired-analysis-v2/physionet2012_mortality/d_gbdt_lr_auroc/mean:4} AUROC, $p=$\rhval{analysis/paired-analysis-v2/physionet2012_mortality/p_gbdt_lr_auroc/mean:4}), with equal Brier score (difference \rhval{analysis/paired-analysis-v2/physionet2012_mortality/d_gbdt_lr_brier/mean:4}, $p=$\rhval{analysis/paired-analysis-v2/physionet2012_mortality/p_gbdt_lr_brier/mean:4}), so its advantage is in ranking rather than in calibration.

\textbf{H2: mask and delta inputs help a GRU over forward-filling (not supported at 48 h).} The mask+delta GRU differs from the forward-fill GRU by \rhval{analysis/paired-analysis-v2/physionet2012_mortality/d_md_ffill_auroc/mean:4} AUROC ($p=$\rhval{analysis/paired-analysis-v2/physionet2012_mortality/p_md_ffill_auroc/mean:4}) and GRU-D by \rhval{analysis/paired-analysis-v2/physionet2012_mortality/d_grud_ffill_auroc/mean:4} ($p=$\rhval{analysis/paired-analysis-v2/physionet2012_mortality/p_grud_ffill_auroc/mean:4}); the AUPRC differences are \rhval{analysis/paired-analysis-v2/physionet2012_mortality/d_md_ffill_auprc/mean:4} and \rhval{analysis/paired-analysis-v2/physionet2012_mortality/d_grud_ffill_auprc/mean:4} (both $p>\rhval{const/alpha:2}$, see Table~\ref{tab:paired}). The effect is therefore indistinguishable from zero here with this number of seeds. The hypothesis is not supported at the full horizon, and holds at shorter horizons (Section~\ref{sec:ablations}).

\textbf{H3: decay adds over mask+delta (not supported).} GRU-D minus the mask+delta GRU is \rhval{analysis/paired-analysis-v2/physionet2012_mortality/d_grud_md_auroc/mean:4} AUROC ($p=$\rhval{analysis/paired-analysis-v2/physionet2012_mortality/p_grud_md_auroc/mean:4}, GRU-D is better in \rhval{analysis/paired-analysis-v2/physionet2012_mortality/w_grud_md_auroc/mean:0} of the five seeds); the decay ablations of Section~\ref{sec:ablations} point the same way.

\textbf{Per-regime breakdown and failure case.} Table~\ref{tab:icu} splits the pooled out-of-fold AUROC by ICU type. GBDT leads in every unit and its margin is largest in the cardiac-surgery recovery unit (CSRU: \rhval{main/gbdt-summary/physionet2012_mortality/auroc_icu2/mean:3} against \rhval{main/gru-d/physionet2012_mortality/auroc_icu2/mean:3} for GRU-D), the unit with the lowest mortality (overall mortality \rhval{data/cohort-statistics/physionet2012_mortality/mortality/mean:3}, CSRU \rhval{data/cohort-statistics/physionet2012_mortality/icu2_mortality/mean:3}) and therefore the fewest positive cases, a regime in which the neural models may be disadvantaged by having the least signal to learn from (we did not test this); the medical ICU is hardest for every system (GBDT \rhval{main/gbdt-summary/physionet2012_mortality/auroc_icu3/mean:3}). Among the recurrent models alone the ordering is not systematic: in the coronary care unit the mask+delta GRU (\rhval{main/gru-mask+delta/physionet2012_mortality/auroc_icu1/mean:3}) beats the forward-fill GRU (\rhval{main/gru-forward-fill/physionet2012_mortality/auroc_icu1/mean:3}) by more than the seed standard deviation, but in the CSRU the ordering reverses. These per-unit numbers come from pooled predictions of about \rhval{const/n_stays} patients split four ways with a few hundred deaths in total, so they are descriptive only.

\begin{table*}[t]\centering\caption{Pooled out-of-fold AUROC by ICU type (CCU coronary care, CSRU cardiac-surgery recovery, MICU medical, SICU surgical), mean $\pm$ sd over seeds.}\label{tab:icu}
\begin{tabular}{lccccc}\toprule
System & CCU & CSRU & MICU & SICU & pooled \\\midrule
LR (summary) & \rhval{main/lr-summary/physionet2012_mortality/auroc_icu1/mean:3} $\pm$ \rhval{main/lr-summary/physionet2012_mortality/auroc_icu1/std:3} & \rhval{main/lr-summary/physionet2012_mortality/auroc_icu2/mean:3} $\pm$ \rhval{main/lr-summary/physionet2012_mortality/auroc_icu2/std:3} & \rhval{main/lr-summary/physionet2012_mortality/auroc_icu3/mean:3} $\pm$ \rhval{main/lr-summary/physionet2012_mortality/auroc_icu3/std:3} & \rhval{main/lr-summary/physionet2012_mortality/auroc_icu4/mean:3} $\pm$ \rhval{main/lr-summary/physionet2012_mortality/auroc_icu4/std:3} & \rhval{main/lr-summary/physionet2012_mortality/auroc_pooled/mean:3} $\pm$ \rhval{main/lr-summary/physionet2012_mortality/auroc_pooled/std:3} \\
GBDT (summary) & \rhval{main/gbdt-summary/physionet2012_mortality/auroc_icu1/mean:3} $\pm$ \rhval{main/gbdt-summary/physionet2012_mortality/auroc_icu1/std:3} & \rhval{main/gbdt-summary/physionet2012_mortality/auroc_icu2/mean:3} $\pm$ \rhval{main/gbdt-summary/physionet2012_mortality/auroc_icu2/std:3} & \rhval{main/gbdt-summary/physionet2012_mortality/auroc_icu3/mean:3} $\pm$ \rhval{main/gbdt-summary/physionet2012_mortality/auroc_icu3/std:3} & \rhval{main/gbdt-summary/physionet2012_mortality/auroc_icu4/mean:3} $\pm$ \rhval{main/gbdt-summary/physionet2012_mortality/auroc_icu4/std:3} & \rhval{main/gbdt-summary/physionet2012_mortality/auroc_pooled/mean:3} $\pm$ \rhval{main/gbdt-summary/physionet2012_mortality/auroc_pooled/std:3} \\
GRU (forward-fill) & \rhval{main/gru-forward-fill/physionet2012_mortality/auroc_icu1/mean:3} $\pm$ \rhval{main/gru-forward-fill/physionet2012_mortality/auroc_icu1/std:3} & \rhval{main/gru-forward-fill/physionet2012_mortality/auroc_icu2/mean:3} $\pm$ \rhval{main/gru-forward-fill/physionet2012_mortality/auroc_icu2/std:3} & \rhval{main/gru-forward-fill/physionet2012_mortality/auroc_icu3/mean:3} $\pm$ \rhval{main/gru-forward-fill/physionet2012_mortality/auroc_icu3/std:3} & \rhval{main/gru-forward-fill/physionet2012_mortality/auroc_icu4/mean:3} $\pm$ \rhval{main/gru-forward-fill/physionet2012_mortality/auroc_icu4/std:3} & \rhval{main/gru-forward-fill/physionet2012_mortality/auroc_pooled/mean:3} $\pm$ \rhval{main/gru-forward-fill/physionet2012_mortality/auroc_pooled/std:3} \\
GRU (mask+delta) & \rhval{main/gru-mask+delta/physionet2012_mortality/auroc_icu1/mean:3} $\pm$ \rhval{main/gru-mask+delta/physionet2012_mortality/auroc_icu1/std:3} & \rhval{main/gru-mask+delta/physionet2012_mortality/auroc_icu2/mean:3} $\pm$ \rhval{main/gru-mask+delta/physionet2012_mortality/auroc_icu2/std:3} & \rhval{main/gru-mask+delta/physionet2012_mortality/auroc_icu3/mean:3} $\pm$ \rhval{main/gru-mask+delta/physionet2012_mortality/auroc_icu3/std:3} & \rhval{main/gru-mask+delta/physionet2012_mortality/auroc_icu4/mean:3} $\pm$ \rhval{main/gru-mask+delta/physionet2012_mortality/auroc_icu4/std:3} & \rhval{main/gru-mask+delta/physionet2012_mortality/auroc_pooled/mean:3} $\pm$ \rhval{main/gru-mask+delta/physionet2012_mortality/auroc_pooled/std:3} \\
GRU-D & \rhval{main/gru-d/physionet2012_mortality/auroc_icu1/mean:3} $\pm$ \rhval{main/gru-d/physionet2012_mortality/auroc_icu1/std:3} & \rhval{main/gru-d/physionet2012_mortality/auroc_icu2/mean:3} $\pm$ \rhval{main/gru-d/physionet2012_mortality/auroc_icu2/std:3} & \rhval{main/gru-d/physionet2012_mortality/auroc_icu3/mean:3} $\pm$ \rhval{main/gru-d/physionet2012_mortality/auroc_icu3/std:3} & \rhval{main/gru-d/physionet2012_mortality/auroc_icu4/mean:3} $\pm$ \rhval{main/gru-d/physionet2012_mortality/auroc_icu4/std:3} & \rhval{main/gru-d/physionet2012_mortality/auroc_pooled/mean:3} $\pm$ \rhval{main/gru-d/physionet2012_mortality/auroc_pooled/std:3} \\
\bottomrule\end{tabular}
\end{table*}

\textbf{Calibration and complexity.} The mean predicted risk of GRU-D (\rhval{main/gru-d/physionet2012_mortality/mean_pred/mean:3}) and of logistic regression (\rhval{main/lr-summary/physionet2012_mortality/mean_pred/mean:3}) is close to the observed mortality (\rhval{data/cohort-statistics/physionet2012_mortality/mortality/mean:3}), whereas gradient boosting systematically under-predicts the average risk (\rhval{main/gbdt-summary/physionet2012_mortality/mean_pred/mean:3}; per-seed range \rhval{main/gbdt-summary/physionet2012_mortality/mean_pred/min:3} to \rhval{main/gbdt-summary/physionet2012_mortality/mean_pred/max:3}). Its Brier score is nevertheless on par with logistic regression, so we make no claim about the role of global bias in the Brier differences and did not compute a calibration slope. The recurrent models stop early (GRU-D after \rhval{main/gru-d/physionet2012_mortality/mean_epochs_or_trees/mean:1} epochs on average, against the cap of \rhval{const/max_epochs}), i.e.\ the inner-validation loss stops improving after few passes over \rhval{const/n_train_fold} training patients; we read this as an overfitting regime but did not test it with, e.g., learning curves.
